\section{Staged Study Design}

\subsection{E0: technical validation}

Before pediatric acquisition, the project verifies Athena-specific raw EEG access, channel metadata, amplitude scaling, sample/timestamp semantics, packet-loss representation, Bluetooth buffering, and synchronization with ARC-ADHD event markers. Representative task and regulation movements are used to characterize artifact burden.

This stage is engineering validation and cannot be presented as evidence for the therapeutic mechanism.

\subsection{E1: observational pediatric substudy}

Children perform the frozen ARC-ADHD task while EEG is recorded. The headband does not alter the intervention, cue selection, task difficulty, recovery criterion, or strategy sequence.

The primary E1 questions are whether synchronized usable EEG can be acquired across the relevant event phases, whether prespecified spectral dynamics are reproducible within child, whether \(T_{\mathrm{EEG}}\) is estimable, and whether neural and behavioral recovery show temporal coupling.

\subsection{E2: incremental-value validation}

Only after E1 measurement gates are met should EEG be tested for incremental information.

A behavioral/context model is compared with the same model plus EEG features under participant-grouped validation. Model selection, preprocessing, feature families, and the primary incremental-value metric must be frozen before evaluation.

Real-time EEG-driven adaptation or neurofeedback belongs to a separate future experiment and is not an automatic next step.
